\chapter{Минимальная source-parent архитектура центральной щели}
\label{ch:v8-baryon-c0-singlet-triplet-central-gap-minimal-source-parent-architecture-gate}

Предыдущий гейт показал, что нормировка не выбирает коэффициент
$\lambda$. Проверим минимальность архитектуры с линейным источником и
положительной квадратичной жёсткостью, не приписывая её параметрам
внутреннее происхождение.

\section{Классификация квадратичных родителей}

Следово-центрированная центральная алгебра одномерна:
\begin{equation}
 \mathcal Z_0=\mathbb R Q,
 \qquad Q=P_{\mathbf3}-\frac34I_4,
 \qquad \Tr Q^2=\frac34.
 \label{eq:v8-baryon-c0-central-gap-source-central-line}
\end{equation}
Общий вещественный многочлен степени не выше двух имеет вид
\begin{equation}
 V(\lambda)=c_0+c_1\lambda+c_2\lambda^2.
 \label{eq:v8-baryon-c0-central-gap-general-quadratic-parent}
\end{equation}
После удаления несущественной константы он имеет единственный ненулевой
глобальный минимум в точности при
\begin{equation}
 c_2>0,
 \qquad c_1\ne0,
 \qquad \lambda_*=-\frac{c_1}{2c_2}.
 \label{eq:v8-baryon-c0-central-gap-quadratic-classification}
\end{equation}
Каждый отдельный моном условия проваливает:
\begin{equation}
 c_1\lambda\ \text{не ограничен снизу при }c_1\ne0,
 \qquad
 c_2\lambda^2\ \text{имеет минимум }\lambda_*=0\text{ при }c_2>0.
 \label{eq:v8-baryon-c0-central-gap-single-monomial-failures}
\end{equation}
Следовательно, степень два и два ненулевых монома являются минимальными.

\section{Источник и жёсткость}

Удобная параметризация классифицированного семейства равна
\begin{equation}
 V_{\rm src}(\lambda)=\frac{m^2}{2}\lambda^2-j\lambda,
 \qquad m^2>0,
 \qquad j\ne0.
 \label{eq:v8-baryon-c0-central-gap-source-parent}
\end{equation}
Она даёт
\begin{equation}
 V'_{\rm src}(\lambda_*)=0,
 \qquad \lambda_*=\frac{j}{m^2},
 \qquad V''_{\rm src}(\lambda_*)=m^2>0.
 \label{eq:v8-baryon-c0-central-gap-source-vacuum-hessian}
\end{equation}
Глобальность и единственность видны из точного завершения квадрата:
\begin{equation}
 V_{\rm src}(\lambda)
 =\frac{m^2}{2}\left(\lambda-\frac{j}{m^2}\right)^2
 -\frac{j^2}{2m^2},
 \qquad V_{\rm src}(\lambda_*)=-\frac{j^2}{2m^2}.
 \label{eq:v8-baryon-c0-central-gap-source-completed-square}
\end{equation}

\section{Операторное поднятие}

Для $h^0=\lambda Q$ тот же родитель записывается без выбора координатного
масштаба на уже нормированной центральной прямой:
\begin{equation}
 \mathcal S_{\rm src}[h^0]
 =\frac{m^2}{2\Tr Q^2}\Tr\bigl((h^0)^2\bigr)
 -\frac{j}{\Tr Q^2}\Tr(Qh^0).
 \label{eq:v8-baryon-c0-central-gap-operator-source-parent}
\end{equation}
Его обратный образ точно равен скалярному потенциалу:
\begin{equation}
 h^0=\lambda Q
 \quad\Longrightarrow\quad
 \mathcal S_{\rm src}[h^0]
 =\frac{m^2}{2}\lambda^2-j\lambda.
 \label{eq:v8-baryon-c0-central-gap-operator-pullback}
\end{equation}
Направление $Q$ вещественно и центрально относительно требуемых действий:
\begin{equation}
 [Q,\rho_{\mathbf1\oplus\mathbf3}(SO(3)_{\rm fam})]=0,
 \qquad [Q,\Gamma_4]=0,
 \qquad \overline Q=Q.
 \label{eq:v8-baryon-c0-central-gap-source-symmetry-compatibility}
\end{equation}
Линейный член нарушает только искусственное отражение координаты
$\lambda\mapsto-\lambda$. Оно не является обязательной внутренней
симметрией, поскольку секторы разных размерностей нельзя переставить
унитарно:
\begin{equation}
 \rank P_{\mathbf1}=1\ne3=\rank P_{\mathbf3}
 \quad\Longrightarrow\quad
 \nexists U:\ UP_{\mathbf1}U^*=P_{\mathbf3}.
 \label{eq:v8-baryon-c0-central-gap-no-required-sign-symmetry}
\end{equation}

\section{Остаточная неидентифицируемость}

Положение вакуума сохраняется на луче параметров
\begin{equation}
 (m^2,j)\longmapsto(am^2,aj),
 \qquad a>0,
 \qquad \frac{aj}{am^2}=\lambda_*.
 \label{eq:v8-baryon-c0-central-gap-source-parameter-orbit}
\end{equation}
При этом гессиан и глубина минимума масштабируются. Точные свидетели
одного вакуума с разной жёсткостью равны
\begin{equation}
 (m^2,j)=(1,1)\Longrightarrow(\lambda_*,V'')=(1,1),
 \qquad
 (m^2,j)=(2,2)\Longrightarrow(\lambda_*,V'')=(1,2).
 \label{eq:v8-baryon-c0-central-gap-source-same-vacuum-witnesses}
\end{equation}
Лишь совместное знание вакуума и физической кривизны $\kappa$ восстанавливает
оба коэффициента:
\begin{equation}
 \lambda_*=\ell,
 \qquad V''(\lambda_*)=\kappa
 \quad\Longleftrightarrow\quad
 m^2=\kappa,
 \qquad j=\kappa\ell.
 \label{eq:v8-baryon-c0-central-gap-source-vacuum-curvature-reconstruction}
\end{equation}

Итоговые реестры разделяют архитектуру и происхождение:
\begin{equation}
 N_{\rm architecture}=7/7,
 \qquad
 N_{\rm parent\ origin}=0/2\quad\{m^2,j\}.
 \label{eq:v8-baryon-c0-central-gap-source-ledgers}
\end{equation}

\begin{theorem}[Минимальность source-parent архитектуры]
Среди вещественных многочленов степени не выше двух единственный класс,
имеющий ограниченный снизу потенциал и единственный ненулевой знаковый
минимум, содержит положительный квадратичный и ненулевой линейный члены.
После удаления константы он эквивалентен $m^2\lambda^2/2-j\lambda$.
Операторное поднятие на $\mathcal Z_0$ сохраняет семейную симметрию,
градуировку и вещественную структуру.
\end{theorem}

\begin{nogocriterion}[Архитектура не является происхождением параметров]
Знание $\lambda_*$ определяет только отношение $j/m^2$. Нельзя объявлять
$j$ и $m^2$ выведенными без независимого вычисления физической кривизны
или эквивалентного двухточечного отклика из одного родительского действия.
\end{nogocriterion}

\begin{protocol}\begin{enumerate}
\item минимальная степень: \textbf{$2$};
\item минимальное число непостоянных мономов: \textbf{$2$};
\item единственный глобальный минимум: \textbf{да при $m^2>0$};
\item ненулевой знак выбран: \textbf{да при $j\ne0$};
\item совместимость с $SO(3)_{\rm fam}$, grading и Real: \textbf{да};
\item условный архитектурный допуск: \textbf{$7/7$};
\item происхождение $m^2,j$: \textbf{$0/2$};
\item следующий гейт: \textbf{parent-origin источника и жёсткости}.
\end{enumerate}\end{protocol}

Машинный сертификат сохранён в
\path{s2t_v8_baryon_c0_singlet_triplet_central_gap_minimal_source_parent_architecture_gate_results.json}.